\src{383}
thương của $G_T$ phẳng trên $T$. Dường như bài toán có nghiệm bằng một lược đồ
$C$ kiểu hữu hạn địa phương trên $S$ khi $X$ riêng trên tiền lược đồ Noether
địa phương $S$ và $G$ còn coherent. Dù sao, chỉ cần giả sử $S$ Noether địa
phương thì hạn chế của $F$ lên các ``$S$-đại số Artin'' là pro-biểu diễn được;
và nếu thêm rằng $X$ riêng trên $S$ và $G$ coherent thì các thành phần địa
phương của vành tôpô tương ứng $O$ là Noether. Dĩ nhiên, ngay cả sau khi chứng
minh sự tồn tại của ``lược đồ Chow'' của $X$ trên $S$, vẫn còn phải tìm một
phân tích thành các tập mở rời nhau $C_i$ (tương ứng với việc cố định các bất
biến liên tục như bậc và chiều của các chu trình đang biến thiên) kiểu hữu hạn
(chứ không chỉ kiểu hữu hạn địa phương) trên $S$, làm rõ quan hệ giữa lược đồ
này với các đa tạp Chow cổ điển, và xác định khi nào một $C_i$ thậm chí xạ ảnh,
hoặc ít nhất tựa xạ ảnh, trên $S$.

\subsection*{3. Các lược đồ Picard}
\addcontentsline{toc}{subsection}{3. Các lược đồ Picard}

Cho $f:X\longrightarrow S$ là một $S$-tiền lược đồ, xét bó nhân
$\mathcal O_X^*$ các đơn vị của bó cấu trúc của $X$, và nhóm
\[
  P(X/S)=H^0\bigl(S,R^1f_*(\mathcal O_X^*)\bigr),
\]
gọi là nhóm Picard tương đối của $X/S$. Như vậy, một phần tử của nhóm này được
xác định bằng cách cho một phủ mở $(U_i)$ của $S$ và một bó khả nghịch $L_i$
trên mỗi $f^{-1}(U_i)$, sao cho với mọi cặp $(i,j)$,
$L_i|f^{-1}(U_i\cap U_j)$ đẳng cấu với
$L_j|f^{-1}(U_i\cap U_j)$, ít nhất là địa phương trên $U_i\cap U_j$ (tức hai bó
này ``tương đương'' theo nghĩa của [3], B, 4). Nếu $X/S$ có một tiết diện thì
$P(X/S)$ không gì khác hơn tập các lớp bó khả nghịch trên $X/S$ modulo
``tương đương'' (loco citato). Bây giờ, với mọi $T$ trên $S$, đặt
\[
  F(T)=P(X_T/T) ;
\]
ta được một hàm tử hiệp biến theo $T$, có thể gọi là hàm tử Picard của $X/S$;
nếu hàm tử này biểu diễn được thì tiền lược đồ trên $S$ biểu diễn nó được gọi
là tiền lược đồ Picard của $X/S$ và ký hiệu $\underline P(X/S)$. Do đó có một
đẳng cấu hàm tử
\[
  \operatorname{Hom}_S\bigl(T,\underline P(X/S)\bigr)
  \xrightarrow{\sim}P(X_T/T).
\]
Phép lập tiền lược đồ Picard tương thích với mở rộng cơ sở; đặc biệt, các tiền
lược đồ Picard của các thớ của $X$ trên $S$ (là các tiền lược đồ trên các
trường thặng dư $K(s)$ của $s\in S$) chính là các thớ của
$\underline P(X/S)$. Dĩ nhiên, vì $P(X_T/T)=F(T)$ là một nhóm giao hoán, các
tiền lược đồ Picard là các tiền lược đồ nhóm. Cũng cần lưu ý rằng các Jacobian
tổng quát của Rosenlicht không gì khác hơn các thành phần liên thông của đơn
vị trong các lược đồ Picard của những đường cong đầy đủ có thể có kỳ dị; điều
này hẳn làm hiển nhiên phần lớn các tính chất của chúng (một khi đã chứng minh
sự tồn tại).

\result{NHẬN XÉT}{} Định nghĩa được chọn ở đây chỉ hợp lý khi mọi điểm của
$Y$ có một lân cận mở $U$ mà trên đó $X$ có một tiết diện. Trong trường hợp
tổng quát, phải sửa nhẹ định nghĩa hàm tử Picard để vẫn có thể thu được một
định lý tồn tại.

Ở đây, các điều kiện hợp lý cho sự tồn tại của một tiền lược đồ Picard là:
$X$ riêng và phẳng trên $S$, $f_*(\mathcal O_X)=\mathcal O_S$, và $X$ có địa
phương một tiết diện trên $S$. Điều kiện này xuất hiện tự nhiên khi áp dụng kỹ
thuật hạ xuống: để loại bỏ các tự đẳng cấu của một bó khả nghịch $L$ trên $X$,
ta trang bị cho nó một tiết diện được đánh dấu trên tiết diện $s$ ([3], B, 4).
Đặc biệt, ta thu được:

